% Level 2, Week 2 -- Functions That Give Back
\lesson{2}{Functions That Give Back}{return hands an answer back; parameters make one function work for many jobs.}
\lessoninfo{Concepts: parameters (R), return (I), docstrings (I), scope (I, light)}{Tools: Thonny, function-machine cards}{CSTA: 2-AP-13, 2-AP-14, 2-AP-19}

\begin{objectives}
  \item The student can explain the difference between printing an answer and returning it, and use a returned value in a new calculation. (L2.G3; L2.W02.1)
  \item The student can write a function with parameters that returns an answer, such as \texttt{damage(attack, shield)}. (L2.G3; L2.W02.2)
  \item The student can call one drawing function several times with different arguments to build a picture. (L2.G7; L2.W02.3)
  \item The student can write a docstring that says what a function gives back, and say why a variable made inside a function is not visible outside it. (L2.G3; L2.W02.4)
\end{objectives}

\begin{materials}
  \item Print (one per student): workbook chapter 2, homework 2.
  \item Function-machine cards (5 minutes): a paper bag or box labelled \textbf{FUNCTION}, about 10 small cards with numbers
        (2, 3, 5, 7, 10 \ldots) and a few with words (\enquote{cat}, \enquote{dragon}), blank cards and a pen.
        Secret rule cards, folded: \enquote{double it}, \enquote{add 3}, \enquote{times 10}, \enquote{count the letters}.
  \item On each Mac: open \texttt{2-intermediate/excercises/week02/ex-game-math/starter.py} in Thonny (\keys{Cmd}+\keys{O}).
  \item A timer, paper and pencils.
\end{materials}

\needspace{16\baselineskip}
\section{At a glance}
Times are a guide: set a timer for each phase and move on when it rings.
\begin{timeplan}
  \timeslot{0--8}{Warm-up}{Unplugged \enquote{Function machines}: guess the secret rule from the cards that go in and come out.}
  \timeslot{8--18}{Concept}{Parameters, arguments and \texttt{return}. Print versus return. Docstrings. What happens inside a function stays inside.}
  \timeslot{18--28}{Guided demo}{Run the Game Math starter, find the \texttt{None} bug together and fix \texttt{xp\_needed()}.}
  \timeslot{28--52}{Hands-on build}{Pairs finish Exercise 2.1 Game Math, then 2.2 Flower Garden. Swap every 8 minutes. Last 5 minutes: garden free-build.}
  \timeslot{52--60}{Wrap-up}{Garden show-and-tell, comment swap, exit ticket, reflection, homework 2.}
\end{timeplan}

\section{Warm-up (0--8 min)}
\begin{unplugged}[Function machines]
You are the machine first. Secretly pick the rule \enquote{double it}. A student hands a number card \textbf{in};
you write the answer on a blank card and hand it \textbf{back out}. After three cards, students guess the rule.
Then a student takes a secret rule card and becomes the machine. Use \enquote{count the letters} last: a \emph{word} goes in, a \emph{number} comes out.
\par\textbf{Point out:} the machine does not shout the answer across the room. It \textbf{hands it back} to the person who asked,
so they can use it (for example, add two answers together). That is exactly what \texttt{return} does.
\end{unplugged}

\section{Concept (8--18 min)}
\begin{teacherscript}[Say]
\enquote{In Level 1 our functions \emph{did} something, like drawing a house. Today our functions also \emph{give back} an answer.
The values we hand in are called \textbf{arguments}; inside the function they land in \textbf{parameters}, which are like labelled slots.
The word \texttt{return} hands the answer back to whoever called the function. Then we can store it, print it, or do maths with it.}
\end{teacherscript}
Write this on paper or show it in Thonny, and trace it with your finger:
\begin{pycode}[Print or return?]
def double_print(n):
    print(n * 2)          # shows 8 on screen, gives back nothing (None)

def double_return(n):
    return n * 2          # gives back 8 to whoever asked

answer = double_return(4) + 1   # 9
oops = double_print(4) + 1      # shows 8, then crashes: None + 1
\end{pycode}

\begin{analogy}[Shouting versus handing over]
Printing is like shouting the answer out of a window: everyone hears it, but nobody can \emph{keep} it.
Returning is like putting the answer in your friend's hand: now they can put it in their pocket (a variable) or
add it to something else. A function that only prints quietly hands back \texttt{None}, Python's word for \enquote{nothing}.
\end{analogy}

\begin{teacherscript}[Say]
\enquote{Two more things. First, a \textbf{docstring} is a sentence in triple quotes right under the \texttt{def} line. It tells
future-you and your friends what the function gives back. Second, variables made inside a function stay inside it, like bowls
in a restaurant kitchen: the diners never see the bowls, only the plate that comes out. If you need a value outside, return it.}
\end{teacherscript}

\section{Guided demo (18--28 min)}
\begin{workedexample}[Finding the None bug in Game Math]
\begin{enumerate}[leftmargin=1.4em]
  \item Open \texttt{ex-game-math/starter.py} and press \menu{Run} (\keys{F5}). The first two lines of output are:
\begin{outputcode}
500
To reach level 5 you need None XP.
\end{outputcode}
  \item Ask: \enquote{Why is 500 on its own line, and why does the sentence say None?} Let students explain: \texttt{xp\_needed} printed
        500 itself, then handed back nothing, so the f-string filled in \texttt{None}.
  \item Change \texttt{print(level * 100)} to \texttt{return level * 100}. Run again: \enquote{you need 500 XP}. Celebrate the fix.
  \item Point at the docstring under \texttt{def xp\_needed(level):}. Read it aloud: \enquote{Return how many experience points you need.}
        Say: \enquote{The docstring promised a return, so the code must keep that promise.}
  \item Show TODO 2 in \texttt{damage()}: ask the group what \texttt{damage(3, 5)} should be. (Not $-2$: a weak goblin cannot heal you! It should be 0.)
        Leave the rest for the pairs.
\end{enumerate}
\end{workedexample}
The finished Game Math, for you:
\pyfile{../excercises/week02/ex-game-math/solution.py}

\section{Hands-on build (28--52 min)}
\begin{pairroles}[Predict, then run]
The \textbf{driver} types. Before each run, the \textbf{navigator} predicts out loud what the function will give back
(for example \enquote{damage(12, 5) should be 7}). Swap at the 8-minute timer.
\textbf{1 student:} you are the navigator and make the predictions; the student types.
\textbf{3 students:} driver, navigator and \textbf{tester}; the tester runs \texttt{check\_my\_work.py} and reads the hints aloud.
\textbf{4 students:} two pairs; after Game Math each pair swaps computers and checks the other pair's docstrings.
\end{pairroles}

\begin{enumerate}[leftmargin=1.4em]
  \item \textbf{Exercise 2.1 Game Math} (\texttt{week02/ex-game-math}, about 10 minutes): TODO 2--5.
  \item \textbf{Exercise 2.2 Flower Garden} (\texttt{week02/ex-flower-garden}, about 10 minutes): make \texttt{petal\_turn()} return
        \texttt{360 / petals} and make \texttt{draw\_flower()} use its \texttt{x}, \texttt{y} and \texttt{petals} parameters.
  \item Check each with \texttt{check\_my\_work.py}.
\end{enumerate}

\begin{bughunt}[Printing instead of returning]
The tell-tale sign is the word \texttt{None} in the output, or the error
\texttt{TypeError: unsupported operand type(s) for +: 'NoneType' and 'int'}.
Ask: \enquote{Which function gave back None? Does it have a \texttt{return} line? Does \emph{every} path through it reach a \texttt{return}?}
Sneaky version: a \texttt{return} only inside an \texttt{if}, so some inputs never reach one.
\end{bughunt}

\begin{freebuild}[Design your garden]
Last 5 minutes: add flowers to the garden by calling \texttt{draw\_flower()} again with new arguments: different places,
petals (try 3 or 20), sizes and colours from the Okabe-Ito list in the file. Fast finishers can draw a row of flowers with
\texttt{for x in range(-200, 201, 100):}.
\end{freebuild}

\section{Wrap-up (52--60 min)}
Show the gardens: each student names one argument they changed.

\begin{safety}[Comments for future-you and friends]
Swap computers for one minute. Read your partner's docstrings and comments. Ask: \enquote{Could someone who has never seen this
code understand what each function gives back?} Agree on three rules: comments explain \emph{why}, they are always kind
(no jokes about people), and they never contain personal information. If an idea came from a friend or a book, a comment can say so:
that is giving credit.
\end{safety}

\begin{formativecheck}[Exit ticket]
What happens when this runs? How would you fix it?
\begin{pycode}
def triple(n):
    print(n * 3)

score = triple(4)
print(score + 1)
\end{pycode}
\textbf{Answer:} it shows \texttt{12}, then crashes with \texttt{TypeError} because \texttt{triple} gave back \texttt{None}.
Fix: \texttt{return n * 3}; then it prints \texttt{13}.
\textbf{Look for:} students who say \enquote{13} straight away have not yet separated showing from giving back. Revisit the analogy next week.
\end{formativecheck}

\begin{reflection}[Say it or write it]
\begin{itemize}[leftmargin=1.2em]
  \item In your own words: what is the difference between \texttt{print} and \texttt{return}?
  \item Which parameter of \texttt{draw\_flower()} was the most fun to change?
\end{itemize}
\end{reflection}
Hand out homework 2 (Potion Shop Calculator).

\section{Notes for the teacher}
\begin{commonmistakes}
  \item \textbf{Print instead of return.} The output shows \texttt{None}. Fix: \texttt{return} the value; print it where the function is called.
  \item \textbf{Calling and throwing away.} \texttt{damage(12, 5)} on its own line works out 7 and loses it. Store it: \texttt{hit = damage(12, 5)}.
  \item \textbf{Arguments in the wrong order.} \texttt{damage(5, 12)} is a different question from \texttt{damage(12, 5)}. Parameters fill up in order.
  \item \textbf{Using an inside variable outside.} \texttt{print(result)} outside \texttt{damage()} gives \texttt{NameError}: \texttt{result} lives only inside.
  \item \textbf{Hard-coded values.} \texttt{draw\_flower()} that still says \texttt{goto(0, 0)} or \texttt{range(4)} ignores its parameters, so every flower looks the same.
\end{commonmistakes}

\begin{differentiation}
\supportpath Do TODO 1 and 3 of Game Math together, then let the pair do TODO 2 and 4. In Flower Garden, do TODO 1 and 2 only
(\texttt{petal\_turn} and \texttt{goto(x, y)}). Keep the \enquote{shouting versus handing over} card on the table.
\extensionpath Write \texttt{heal(health, potion)} that returns the new health but never more than 100. Explain the default
value in \texttt{size=60}: what happens if you leave it out of the call? Add a docstring to every function.
\end{differentiation}

\begin{stretchchallenge}[A whole meadow]
Write \texttt{draw\_row(t, y, how\_many)} that draws \texttt{how\_many} flowers in a row at height \texttt{y} by calling
\texttt{draw\_flower()} in a loop. Then call it three times for a meadow. Bonus: a function \texttt{petals\_for(row)} that
\emph{returns} a different number of petals for each row.
\end{stretchchallenge}

\begin{teachernote}[If you are short on time]
Shorten the function-machine game to two rounds. Do the Flower Garden as a whole-group demo (you drive, students navigate)
and keep Game Math as the pair task. Never skip the print-versus-return bug: it is the key idea of the week.
\end{teachernote}
